\documentclass[10pt,a4paper]{amsart}
\usepackage{amsmath,amssymb,amsthm,graphicx}
% Retained minimal preamble for the Stacks Project chapter more-morphisms.tex (More on Morphisms).
% Verbatim excerpt of the chapter preamble: Xy-pic setup (native preamble lines 14-18) and the native
% theorem declarations with the mathematical macros (native preamble lines 154-217). Copied mechanically;
% nothing added, removed or redefined. The wrapper class and the amsmath/amssymb/amsthm/graphicx base
% packages are supplied by the assembler, so the class/externaldocument part of the native preamble is omitted.
% For commutative diagrams we use Xy-pic
\usepackage[all]{xy}

% We use 2cell for 2-commutative diagrams.
\xyoption{2cell}
% Theorem environments.
%
\theoremstyle{plain}
\newtheorem{theorem}[subsection]{Theorem}
\newtheorem{proposition}[subsection]{Proposition}
\newtheorem{lemma}[subsection]{Lemma}

\theoremstyle{definition}
\newtheorem{definition}[subsection]{Definition}
\newtheorem{example}[subsection]{Example}
\newtheorem{exercise}[subsection]{Exercise}
\newtheorem{situation}[subsection]{Situation}

\theoremstyle{remark}
\newtheorem{remark}[subsection]{Remark}
\newtheorem{remarks}[subsection]{Remarks}

\numberwithin{equation}{subsection}

% Macros
%
\def\lim{\mathop{\mathrm{lim}}\nolimits}
\def\colim{\mathop{\mathrm{colim}}\nolimits}
\def\Spec{\mathop{\mathrm{Spec}}}
\def\Hom{\mathop{\mathrm{Hom}}\nolimits}
\def\Ext{\mathop{\mathrm{Ext}}\nolimits}
\def\SheafHom{\mathop{\mathcal{H}\!\mathit{om}}\nolimits}
\def\SheafExt{\mathop{\mathcal{E}\!\mathit{xt}}\nolimits}
\def\Sch{\mathit{Sch}}
\def\Mor{\mathop{\mathrm{Mor}}\nolimits}
\def\Ob{\mathop{\mathrm{Ob}}\nolimits}
\def\Sh{\mathop{\mathit{Sh}}\nolimits}
\def\NL{\mathop{N\!L}\nolimits}
\def\CH{\mathop{\mathrm{CH}}\nolimits}
\def\proetale{{pro\text{-}\acute{e}tale}}
\def\etale{{\acute{e}tale}}
\def\QCoh{\mathit{QCoh}}
\def\Ker{\mathop{\mathrm{Ker}}}
\def\Im{\mathop{\mathrm{Im}}}
\def\Coker{\mathop{\mathrm{Coker}}}
\def\Coim{\mathop{\mathrm{Coim}}}

% Boxtimes
%
\DeclareMathSymbol{\boxtimes}{\mathbin}{AMSa}{"02}

%
% Macros for moduli stacks/spaces
%
\def\QCohstack{\mathcal{QC}\!\mathit{oh}}
\def\Cohstack{\mathcal{C}\!\mathit{oh}}
\def\Spacesstack{\mathcal{S}\!\mathit{paces}}
\def\Quotfunctor{\mathrm{Quot}}
\def\Hilbfunctor{\mathrm{Hilb}}
\def\Curvesstack{\mathcal{C}\!\mathit{urves}}
\def\Polarizedstack{\mathcal{P}\!\mathit{olarized}}
\def\Complexesstack{\mathcal{C}\!\mathit{omplexes}}
% \Pic is the operator that assigns to X its picard group, usage \Pic(X)
% \Picardstack_{X/B} denotes the Picard stack of X over B
% \Picardfunctor_{X/B} denotes the Picard functor of X over B
\def\Pic{\mathop{\mathrm{Pic}}\nolimits}
\def\Picardstack{\mathcal{P}\!\mathit{ic}}
\def\Picardfunctor{\mathrm{Pic}}
\def\Deformationcategory{\mathcal{D}\!\mathit{ef}}
\newenvironment{reference}{}{}
\newenvironment{slogan}{}{}
\newenvironment{history}{}{}
\usepackage{cleveref}
\crefname{originallemma}{original lemma}{original lemmas}
\Crefname{originallemma}{Original lemma}{Original lemmas}
\makeatletter
\newlabel{lemma-pushout-along-thickening}{{lemma-pushout-along-thickening}{source}{}{sourcepointer-1}{}}
\newlabel{lemma-pushout-along-thickening@cref}{{[originallemma][0][]lemma-pushout-along-thickening}{[1][0][]source}{}{}{}}
\makeatother
\crefname{originallemma}{original lemma}{original lemmas}
\Crefname{originallemma}{Original lemma}{Original lemmas}
\makeatletter
\newlabel{more-algebra-lemma-module-over-fibre-product}{{more-algebra-lemma-module-over-fibre-product}{source}{}{sourcepointer-2}{}}
\newlabel{more-algebra-lemma-module-over-fibre-product@cref}{{[originallemma][0][]more-algebra-lemma-module-over-fibre-product}{[1][0][]source}{}{}{}}
\makeatother
\crefname{originallemma}{original lemma}{original lemmas}
\Crefname{originallemma}{Original lemma}{Original lemmas}
\makeatletter
\newlabel{morphisms-lemma-base-change-affine}{{morphisms-lemma-base-change-affine}{source}{}{sourcepointer-3}{}}
\newlabel{morphisms-lemma-base-change-affine@cref}{{[originallemma][0][]morphisms-lemma-base-change-affine}{[1][0][]source}{}{}{}}
\makeatother
\title{Schemes over a pushout sit in a category equivalence with pairs of schemes glued along the closed subscheme}
\author{The Stacks Project Authors}
\date{stacks-project:more-morphisms.tex:local-snapshot-20261011; GNU FDL 1.2 or later}
\begin{document}
\maketitle
\noindent\textit{Editorial note.} Statement span is lines 4276-4293 of sources/stacks-more-morphisms/pinned-original-source.tex; proof span is lines 4295-4346, the complete author proof (50 body lines).\par\medskip
\noindent\textit{Editorial note.} Navigation stubs: 3 cited-but-not-included labels are supplied as non-load-bearing pointers with their pinned tags where the snapshot tag index has them.\par\medskip
\noindent\textit{Editorial note.} extra\_packages is empty on purpose: the retained preamble already loads xy with options and a second declaration clashes in the assembler wrapper.\par\medskip
\noindent\textit{Editorial note.} Text is retained verbatim, including the author's own typos and placeholders.\par\medskip
\noindent\textit{Allowed standard prerequisites (editorial).}\begin{itemize}
\item Explicit prerequisite boundary. Pushouts of schemes along closed immersions, gluing of schemes, and the universal property of a pushout square; the equivalence is retained.
\end{itemize}\medskip
\ifdefined\hypertarget\hypertarget{sourcepointer-1}{}\fi
\noindent\textit{Source reference.} Original lemma lemma-pushout-along-thickening: sources/stacks-more-morphisms/pinned-original-source.tex, source lines 4207--4227. Cited by the retained statement or proof but not part of the collected spans; supplied only as a navigation pointer. The argument does not depend on its content.\par\medskip
\ifdefined\hypertarget\hypertarget{sourcepointer-2}{}\fi
\noindent\textit{Source reference.} Original lemma more-algebra-lemma-module-over-fibre-product: sources/stacks-more-morphisms/pinned-original-source.tex, source lines 4346--4346. Cited by the retained statement or proof but not part of the collected spans; supplied only as a navigation pointer. The argument does not depend on its content.\par\medskip
\ifdefined\hypertarget\hypertarget{sourcepointer-3}{}\fi
\noindent\textit{Source reference.} Original lemma morphisms-lemma-base-change-affine: sources/stacks-more-morphisms/pinned-original-source.tex, source lines 4346--4346. Cited by the retained statement or proof but not part of the collected spans; supplied only as a navigation pointer. The argument does not depend on its content.\par\medskip
\begin{lemma}
\label{lemma-equivalence-categories-schemes-over-pushout}
Let $X \to X'$ be a thickening of schemes and let $X \to Y$ be an
affine morphism of schemes. Let $Y' = Y \amalg_X X'$ be the pushout
(see Lemma \ref{lemma-pushout-along-thickening}). Base change gives
a functor
$$
F : (\Sch/Y') \longrightarrow (\Sch/Y) \times_{(\Sch/X)} (\Sch/X')
$$
given by $V' \longmapsto (V' \times_{Y'} Y, V' \times_{Y'} X', 1)$
which has a left adjoint
$$
G : (\Sch/Y) \times_{(\Sch/X)} (\Sch/X') \longrightarrow (\Sch/Y')
$$
which sends the triple $(V, U', \varphi)$ to the pushout
$V \amalg_{(V \times_Y X)} U'$. Finally, $F \circ G$ is isomorphic to the
identity functor.
\end{lemma}

\begin{proof}
Let $(V, U', \varphi)$ be an object of the fibre product category.
Set $U = U' \times_{X'} X$. Note that $U \to U'$ is a thickening.
Since $\varphi : V \times_Y X \to U' \times_{X'} X = U$ is an isomorphism
we have a morphism $U \to V$ over $X \to Y$ which identifies $U$ with
the fibre product $X \times_Y V$. In particular $U \to V$ is affine, see
Morphisms, Lemma \ref{morphisms-lemma-base-change-affine}.
Hence we can apply Lemma \ref{lemma-pushout-along-thickening}
to get a pushout $V' = V \amalg_U U'$. Denote $V' \to Y'$ the morphism
we obtain in virtue of the fact that $V'$ is a pushout and because
we are given morphisms $V \to Y$ and $U' \to X'$ agreeing on $U$
as morphisms into $Y'$. Setting $G(V, U', \varphi) = V'$
gives the functor $G$.

\medskip\noindent
Let us prove that $G$ is a left adjoint to $F$. Let $Z$ be a scheme
over $Y'$. We have to show that
$$
\Mor(V', Z) = \Mor((V, U', \varphi), F(Z))
$$
where the morphism sets are taking in their respective categories.
Let $g' : V' \to Z$ be a morphism. Denote $\tilde g$, resp.\ $\tilde f'$
the composition of $g'$ with the morphism $V \to V'$, resp.\ $U' \to V'$.
Base change $\tilde g$, resp.\ $\tilde f'$ by $Y \to Y'$, resp.\ $X' \to Y'$
to get a morphism $g : V \to Z \times_{Y'} Y$,
resp.\ $f' : U' \to Z \times_{Y'} X'$. Then $(g, f')$ is an element
of the right hand side of the equation above (details omitted).
Conversely, suppose that $(g, f') : (V, U', \varphi) \to F(Z)$ is an
element of the right hand side.
We may consider the composition $\tilde g : V \to Z$,
resp.\ $\tilde f' : U' \to Z$ of $g$, resp.\ $f$ by
$Z \times_{Y'} X' \to Z$, resp.\ $Z \times_{Y'} Y \to Z$.
Then $\tilde g$ and $\tilde f'$ agree as morphism from $U$ to $Z$.
By the universal property of pushout, we obtain a morphism
$g' : V' \to Z$, i.e., an element of the left hand side.
We omit the verification that these constructions are mutually inverse.

\medskip\noindent
To prove that $F \circ G$ is isomorphic to the identity we have to
show that the adjunction mapping
$(V, U', \varphi) \to F(G(V, U', \varphi))$ is an isomorphism.
To do this we may work affine locally. Say $X = \Spec(A)$, $X' = \Spec(A')$,
and $Y = \Spec(B)$. Then $A' \to A$ and $B \to A$ are ring maps as in
More on Algebra, Lemma \ref{more-algebra-lemma-module-over-fibre-product}
and $Y' = \Spec(B')$ with $B' = B \times_A A'$. Next, suppose that
$V = \Spec(D)$, $U' = \Spec(C')$ and $\varphi$ is given by an
$A$-algebra isomorphism $D \otimes_B A \to C' \otimes_{A'} A = C'/IC'$.
Set $D' = D \times_{C'/IC'} C'$. In this case the statement we have to
prove is that $D' \otimes_{B'} B \cong D$ and $D' \otimes_{B'} A' \cong C'$.
This is a special case of More on Algebra, Lemma
\ref{more-algebra-lemma-module-over-fibre-product}.
\end{proof}

\end{document}
